% !TeX program = xelatex
% !BIB program = biber
\documentclass{tongji-thesis}
\usepackage[backend=biber,style=gb7714-2015]{biblatex}
\addbibresource{references.bib}

\Title{xxxxxxxxxxxxxxxxxxxxxx}
\StudentID{2534xxx}
\StudentName{xxx}
\DisciplineCategory{工学}
\FirstDiscipline{软件工程}
\ResearchField{软件与人工智能}
\College{计算机科学与技术学院}
\JointUnit{}
\Supervisor{xx}
\ProposalDate{2026-05-18}
\ProposalYear{2026}
\ProposalMonth{5}
\ProposalDay{18}

\begin{document}

\makeproposalcover
\proposalrequirement

\section{绪论}

\subsection{xxx}

\subsubsection{xxxx}

\paragraph{xxx}

随着\textbf{深度学习}模型规模不断扩大，\cite{knuth1984texbook}。

\begin{align}
  L &= L_{\mathrm{seg}} + \lambda L_{\mathrm{reg}} \\
    &= -\sum_i y_i \log \hat{y}_i + \lambda \|\theta\|_2^2
\end{align}

\begin{algorithm}[htbp]
\caption{模型训练流程}
\begin{algorithmic}[1]
\Require 训练集 $D$，学习率 $\eta$，最大轮数 $T$
\Ensure 训练后的模型参数 $\theta$
\State 初始化模型参数 $\theta$
\For{$t = 1$ to $T$}
  \State 从 $D$ 中采样一个小批量数据
  \State 计算损失函数 $L$
  \State 根据梯度更新参数 $\theta \leftarrow \theta - \eta \nabla_\theta L$
\EndFor
\State \Return $\theta$
\end{algorithmic}
\end{algorithm}

\newpage
\printbibliography

\newpage
\begin{researchprogresstable}
\progressdate{2025.03}{2025.04} & 查阅文献，完成研究背景与国内外研究现状分析。 & 完成文献综述与选题依据。 \\
\progressdate{2025.05}{2025.06} & 设计研究方案，确定技术路线和实验数据集。 & 完成总体方案设计。 \\
\progressdate{2025.07}{2025.08} & 搭建模型，完成初步实验。 & 获得初步实验结果。 \\
\progressdate{2025.09}{2025.10} & 优化模型结构，进行对比实验与消融实验。 & 完成主要实验分析。 \\
\progressdate{2025.11}{2025.12} & 整理实验结果，撰写论文初稿。 & 完成论文初稿。 \\
\progressdate{2026.01}{2026.02} & 修改论文，完善格式，准备答辩材料。 & 完成论文定稿。 \\
\end{researchprogresstable}

\end{document}
